% ==============================================================================
% SLIDE 7: GEOLOGICAL STRATIGRAPHY & VOLCANIC HAZARDS
% ==============================================================================
\begin{frame}[label=geology]{Caldera Geology: Anisotropic Volcanic Stratigraphy}
  \begin{columns}[T, onlytextwidth]
    % Left Column: Geological Diagram
    \begin{column}{0.48\textwidth}
      \slidephoto{Santorini-volcano.JPG}{Layered stratigraphy of the 300\,m caldera rim: Dense lavas interbedded with friable pumice and scoria}
    \end{column}

    % Right Column: Bullet Points
    \begin{column}{0.48\textwidth}
      \textbf{\color{aegeanblue}Geomechanical Hazard Profile:}
      \vspace{0.1em}
      \begin{itemize}
          \setlength{\itemsep}{0.08em}
        \item \textbf{Non-Monolithic Mass:} Alternating hard dacite/andesite lavas over fragile, unconsolidated pumice and scoria (\textit{tephra}).
        \item \textbf{Rotational Slip Planes:} Weathering creates shear surfaces; shallow footings risk shearing off the cliff edge under loads.
        \item \textbf{Seismic Swarms:} 2011--2012 caldera inflation and tremors from Kolumbo tested the geotechnical limits of cliff hotels.
        \item \textbf{Modern Code Resilience:} Properties built to post-1956 seismic standards endured thousands of tremors without collapse.
      \end{itemize}
    \end{column}
  \end{columns}
\end{frame}
